% Revision de la literatura
\section{Related work}
\label{sec:literature}

Three-dimensional packing problems with practical constraints have received much attention; we
refer to \citet{Bortfeldt2013} for a review of the constraints considered in container loading,
to \citet{Zhao2016} for a comparison of container loading algorithms, and to
\citet{Brandt2019} for a definition of the air cargo load planning problem and a review of its
subproblems. We focus on the works closest to ours: air cargo palletization, packing with
weight distribution constraints, and the heuristics on which our algorithm is based.

\paragraph{Air cargo palletization.}
\citet{Chan2006} developed a decision support system for loading air cargo pallets in a real
case study. \citet{Paquay2016} formulated the problem considered in this paper as a mixed
integer program, with non-cuboid ULDs, weight capacities, orientations, fragile boxes, support
and CoG constraints, and solved small instances with a commercial solver. To solve larger
instances, \citet{Paquay2018a} proposed three MIP-based constructive heuristics (relax-and-fix,
insert-and-fix and fractional relax-and-fix), and \citet{Paquay2018b} a tailored two-phase
constructive heuristic: the first phase packs the boxes into identical bins with an extension of
the extreme points, and the second one deals with the several ULD types; if necessary, the weight distribution of
the best loading is then improved in a post-processing step with shift and jump operators that
move or repack boxes. \citet{Hessler2025}
extended the problem with the ULD edge, the substructure and the padding material, and proposed a
randomized greedy search (RGS). Each ULD is loaded by repeatedly calling an insertion heuristic
with several sorting criteria of the boxes and a random perturbation of the order; the
insertion heuristic places each box at the first feasible extreme point
\citep{Crainic2008}, uses an extended set of extreme points that can be moved to fill the space
under the bevelled sides of the ULD, and accelerates the feasibility checks with a grid of
cells. The best loading of a ULD is selected with a score that combines the volume utilization
and the weight balance, and the ULDs are loaded sequentially, first those that are needed for
the boxes that fit in the fewest ULD types. Their algorithm improves the results of
\citet{Paquay2018b} by about 30\% in a fraction of a second per instance. A related but
different problem is the loading of the ULDs into the aircraft, which determines the position of
each ULD to balance the aircraft, sometimes with several pickup and delivery points
\citep{Lurkin2015,Mesquita2026}.

\paragraph{Weight distribution.}
Constraints on the CoG of the load arise in container loading \citep{Davies1999}, in pallet
loading and in air cargo \citep{Chan2006,Paquay2016}, and they have also been studied in the
three-dimensional knapsack problem \citep{Baldi2012} and in the load-balanced bin packing
problem \citep{Trivella2016}. Many heuristics handle them in a post-processing step, for
instance by reordering and mirroring the walls of the load after the packing has been built
\citep{Davies1999} or by shifting and repacking boxes \citep{Paquay2018b}, or with a score that trades the weight balance against the volume
utilization \citep{Hessler2025}. \citet{Paquay2018b} also consider the CoG during the
construction, with a merit function that selects the extreme point that keeps the CoG closest to
the center of the ULD. In our algorithm the CoG is taken into account in the evaluation of every
candidate block and in the value of the partial solutions that guides the beam search, and the
load is also translated at the end.

\paragraph{Heuristics for three-dimensional packing.}
Most heuristics for three-dimensional bin packing build the packing by placing one box at a
time at a candidate position: the corner points of \citet{Martello2000}, the extreme points of
\citet{Crainic2008}, or the maximal free spaces of \citet{Parreno2008}. \citet{Parreno2008} and
\citet{Parreno2010} embedded a maximal-space constructive heuristic in GRASP and variable
neighbourhood descent, and \citet{AlvarezValdes2013} extended it with path relinking to the
multiple bin size bin packing problem. \citet{Gajda2022} proposed a randomized constructive
heuristic for a real container loading problem with weight distribution along the truck axles.
For the single container loading problem, the best results have been obtained by tree search
algorithms that place \emph{blocks} of boxes instead of single boxes:
\citet{Fanslau2010} combined blocks of boxes with a partition-controlled tree search,
\citet{Zhu2012} proposed an iterative-doubling greedy lookahead algorithm, and
\citet{Araya2014} a beam search in which each node is evaluated by a greedy completion of its
partial solution (BSG). \citet{Araya2017} introduced the VCS function to evaluate the blocks,
which combines the volume of the block, its contact surface with the placed boxes and the
container, the volume that the block makes unusable in the free space, and the number of boxes,
\citet{Araya2020} extended BSG to a biobjective container loading problem that also
maximizes the profit of the loaded boxes, and \citet{Guesser2026} improved its implementation
with cache-friendly data structures, the memoization of the greedy completions and a two-stage
evaluation function, obtaining the best results to date on the classical instances.
\citet{Hessler2025} also tested the RGS on the container loading instances of
\citet{Bischoff1995} and reported that it lags behind BSG with VCS \citep{Araya2017}, which they
attribute to the much larger number of boxes per container in those instances. To the best of
our knowledge, block-based tree search has not been applied to air cargo palletization, where
the boxes are highly heterogeneous, the ULDs are not cuboids, and the support, fragility and CoG
constraints restrict the blocks that can be built and where they can be placed. The main purpose
of this paper is to show how to handle these features and that, with them, block-based tree
search also outperforms the state of the art on this problem.

Finally, the reduction of the volume and the insertion of the unloaded boxes in our multi-ULD
framework are related to the local search heuristics that try to empty one bin at a time
\citep{Faroe2003}, to the repacking of the contents of large bins into smaller ones in
variable-sized bin packing \citep{Ortmann2010}, and to the ruin and recreate principle
\citep{Schrimpf2000}.
